\section*{अध्याय \ref{ch:b1:ocean-circulation} --- महासागर का परिसंचरण}
% values: calc/ch29.py

\begin{solution}{exo:b1:ocean-circulation:1}
(क)~\qty{4}{m/s} पर (पहला सूत्र): $1000\,C_D = 0.29 + 0.775 + 0.481 = 1.546$,
$T = 1.3 \times \num{1.546e-3} \times 16 = \qty{0.0322}{Pa}$। \qty{7}{m/s} पर (दूसरा सूत्र): $1000\,C_D = 0.60 + 0.497 = 1.097$, $T = \qty{0.0699}{Pa}$। \qty{15}{m/s} पर: $1000\,C_D = 1.665$, $T = \qty{0.487}{Pa}$। (ख)~$0.487/0.0699 =
6.97$, $(15/7)^2 = 4.59$ के विरुद्ध: खिंचाव-गुणांक स्वयं पवन के साथ बढ़ता है, क्योंकि प्रबल पवन खुरदरी लहरें उठाती है जो हवा को बेहतर पकड़ती हैं।
\end{solution}

\begin{solution}{exo:b1:ocean-circulation:2}
(क)~उत्तरी गोलार्ध, पूर्व की ओर बहती पवन के दाईं ओर \ang{90}: दक्षिण की ओर। (ख)~पवन दक्षिण-पश्चिम की ओर बहती है; उसके दाईं ओर उत्तर-पश्चिम है। (ग)~दक्षिणी गोलार्ध, पूर्व की ओर बहती पवन के बाईं ओर \ang{90}: उत्तर की ओर। (घ)~पवन उत्तर की ओर बहती है; उसके बाईं ओर पश्चिम है, जो पश्चिमी तट पर अपतटीय है। स्थिति (घ) में परिवहन अपतटीय इंगित करता है: सतही जल नीचे से बदला जाता है, और ठंडा जल तट के साथ उत्सर्पित होता है।
\end{solution}

\begin{solution}{exo:b1:ocean-circulation:3}
(क)~$31.7 + 3.5 - 17.9 = \qty{17.3}{\sverdrup}$, प्रकाशित \qty{17.1}{\sverdrup} के विरुद्ध: तीनों घटकों के माध्य पूर्णांकन और अनिश्चितताओं ($\pm\qty{3.2}{\sverdrup}$) के भीतर अधिक्रमण के माध्य में जुड़ते हैं। (ख)~\qty{1100}{m} के ऊपर: $16.6 + 0.4 = \qty{17.0}{\sverdrup}$ उत्तर की ओर; नीचे: $12.0 + 5.9 = \qty{17.9}{\sverdrup}$ दक्षिण की ओर। दोनों लगभग संतुलित हैं, जैसा होना चाहिए जब खंड के उत्तर में जल जमा न हो; छोटा अंतर अनिश्चितताओं के भीतर है। (ग)~$\num{17.1e6} \times \num{3.156e7} = \qty{5.40e14}{m^3}$, यानी प्रति वर्ष \qty{5.40e5}{km^3}; $\num{1.3324e9}/(\num{5.40e5}) \approx 2470$ वर्ष।
\end{solution}

\begin{solution}{exo:b1:ocean-circulation:4}
$D_E = 7.6 \times 10/\sqrt{\sin\varphi}$ और $V_0 = 0.0127 \times
10/\sqrt{\sin\varphi}$। \ang{15} पर: \qty{149}{m} और \qty{0.250}{m/s}
(पवन का \qty{2.5}{\percent}); \ang{45} पर: \qty{90}{m} और \qty{0.151}{m/s}
(\qty{1.5}{\percent}); \ang{60} पर: \qty{82}{m} और \qty{0.136}{m/s}
(\qty{1.4}{\percent})। भूमध्य रेखा के निकट $\sin\varphi \to 0$ और दोनों सूत्र अपसरित होते हैं; अधिक मूलभूत रूप से $f \to 0$ है, इसलिए घर्षण को संतुलित करने वाला कोरिऑलिस बल शून्य हो जाता है, और संतुलन को स्थापित होने में लगने वाला समय $1/f$ उस अवधि से लंबा हो जाता है जब तक पवन स्थिर रहती है। एकमान ने अपना सूत्र $|\varphi| \geq
\ang{10}$ तक सीमित रखा।
\end{solution}

\begin{solution}{exo:b1:ocean-circulation:5}
(क)~$f = 2\Omega\sin\ang{11} = \qty{2.783e-5}{s^{-1}}$; $M = 0.07/(\num{2.783e-5}) =
\qty{2515}{kg.m^{-1}.s^{-1}}$, पश्चिम-वाही पवन के दाईं ओर \ang{90}: उत्तर की ओर। (ख)~$Q = (2515/1027) \times \num{5e6} = \qty{1.22e7}{m^3/s}$,
\qty{12.2}{\sverdrup}। (ग)~आकलन प्रत्यक्ष मापन की अनिश्चितता ($12.0 \pm 5.5$) के भीतर है और दीर्घकालीन पवनों से मान ($13.5 \pm 0.3$) के निकट: प्रमेय प्रेक्षित परिवहन की व्याख्या करता है।
\end{solution}

\begin{solution}{exo:b1:ocean-circulation:6}
(क)~$f(\ang{2}) = \qty{5.09e-6}{s^{-1}}$; $M = 0.05/(\num{5.09e-6}) =
\qty{9824}{kg.m^{-1}.s^{-1}}$, \qty{9.57}{m^2/s}: \ang{2} उ पर उत्तर की ओर (पश्चिम-वाही पवन के दाईं ओर), \ang{2} द पर दक्षिण की ओर (बाईं ओर)। (ख)~\qty{1}{m} लंबी पट्टी पर \omterm{def:b1:spherical-astronomy:geographic}{देशांतर} के अनुदिश, $2 \times 9.57 = \qty{19.1}{m^3/s}$ दो अक्षांशों के आर-पार बाहर जाते हैं, जो \qty{444.8}{km} दूर हैं: $w = 19.1/(\num{4.448e5}) = \qty{4.30e-5}{m/s}$, लगभग \qty{3.7}{m} प्रति दिन। (ग)~भूमध्य रेखा पर $f = 0$ है और $|\vec M| = T/|f|$ अनंत है: वहाँ स्थायी एकमान संतुलन नहीं है, और ऊपर की संख्याएँ, चुने गए \omterm{def:b1:spherical-astronomy:geographic}{अक्षांश} के प्रति बहुत संवेदनशील, केवल एक कोटि की हैं। जो बचता है वह चिह्न है: दोनों ओर परिवहन भूमध्य रेखा से दूर इंगित करता है, इसलिए सतही जल अपसरित होता है और उसके साथ ठंडा जल ऊपर उठता है।
\end{solution}

\begin{solution}{exo:b1:ocean-circulation:7}
(क)~$f(\ang{15}) = \qty{3.775e-5}{s^{-1}}$: $M = 0.06/(\num{3.775e-5}) =
\qty{1590}{kg.m^{-1}.s^{-1}}$ उत्तर की ओर; $f(\ang{40}) =
\qty{9.375e-5}{s^{-1}}$: $M = \qty{1280}{kg.m^{-1}.s^{-1}}$ दक्षिण की ओर।
(ख)~\qty{2780}{km} दूर दो अक्षांशों के बीच, $\partial M_y/\partial y =
-(1590 + 1280)/(\num{2.78e6}) = \qty{-1.03e-3}{kg.m^{-2}.s^{-1}}$: \cref{prop:b1:ocean-circulation:pumping} से, $w_E = \qty{-1.005e-6}{m/s}$, प्रति वर्ष \qty{31.7}{m} का नीचे की ओर वेग। (ग)~$\num{1.005e-6} \times \num{2.78e6} \times
\num{5e6} = \qty{1.40e7}{m^3/s}$, \qty{14.0}{\sverdrup} नीचे पंप होते हैं।
\end{solution}

\begin{solution}{exo:b1:ocean-circulation:8}
(क)~$H = 1027 \times 3992 \times \num{17e6} \times 12 = \qty{8.36e14}{W}$,
\qty{0.84}{PW}। (ख)~यह मापे गए \qty{1.2}{PW} का \qty{70}{\percent} है: वास्तविक ऊपरी शाखा सतह के निकट \qty{15}{\celsius} के माध्य से गरम है, और \omterm{def:b1:ocean-circulation:gyre}{गायर} भी ऊष्मा ले जाता है। (ग)~$1027 \times 3992 \times \num{17e6} \times 3 =
\qty{2.1e14}{W}$, \qty{0.21}{PW}, (क) का एक-चौथाई। केवल वह जल जो लौटने से पहले ठंडा हो जाता है बहुत ऊष्मा ले जाता है: अटलांटिक का ऊष्मा परिवहन गहरे जल के निर्माण पर निर्भर है।
\end{solution}

\begin{solution}{exo:b1:ocean-circulation:9}
माना $x$, $y$ और $1 - x - y$ AAIW, NADW और AABW के अंश हैं। लवण और ऊष्मा संतुलन, AABW से मापे गए:
\[
  -0.52\,x + 0.33\,y = 34.731 - 34.67 = 0.061,\qquad
  2.9\,x + 3.6\,y = 2.78 - 0.4 = 2.38 .
\]
सारणिक $-0.52 \times 3.6 - 0.33 \times 2.9 = -2.829$ है, और क्रैमर का नियम
देता है $x = (0.061 \times 3.6 - 0.33 \times 2.38)/(-2.829) = 0.20$ और $y = (-0.52
\times 2.38 - 2.9 \times 0.061)/(-2.829) = 0.50$: नमूना \qty{20}{\percent} AAIW, \qty{50}{\percent} NADW और \qty{30}{\percent} AABW है, तीनों जलों के त्रिभुज के भीतर एक बिंदु।
\end{solution}

\begin{solution}{exo:b1:ocean-circulation:10}
(क)~$W = u + iv = V_0e^{i\pi/4}e^{(1+i)az}$ के साथ, जिसके वास्तविक और काल्पनिक भाग
प्रस्ताव के सूत्र हैं, $W'' = (1+i)^2a^2W = 2ia^2W$। दो समीकरण
मिलकर $A_z(u'' + iv'') = -fv + ifu = if(u + iv)$ बनाते हैं, यानी $A_zW'' = ifW$, जो सत्य है क्योंकि $2ia^2A_z = if$ है, $a^2 = f/(2A_z)$ के लिए। (ख)~अब $W'(0) =$
$(1+i)aV_0e^{i\pi/4} =$ $\sqrt2\,aV_0\,e^{i\pi/2} =$ $i\sqrt2\,aV_0$: सतही प्रतिबल $\rho A_zW'(0)$ उत्तर की ओर इंगित करता है, परिमाण
\[ \rho A_z\sqrt2\,aV_0 = \rho A_z\sqrt2\sqrt{f/(2A_z)}\;T/(\rho\sqrt{fA_z}) = T \]
(ग)~अब $\int_{-\infty}^0 W\,\mathrm dz =
V_0e^{i\pi/4}/((1+i)a) = V_0/(\sqrt2\,a)$, एक वास्तविक संख्या: परिवहन का कोई उत्तरी भाग नहीं है, और उसका पूर्वी भाग $\rho V_0/(\sqrt2\,a) =
T/\sqrt{fA_z}\cdot\sqrt{A_z/f} = T/f$ है। यह प्रमेय है: परिमाण $T/f$, उत्तर की ओर बहती पवन के दाईं ओर \ang{90}।
\end{solution}

\begin{solution}{exo:b1:ocean-circulation:11}
(क)~उत्तर की ओर
\[ \mathrm dy = a_\oplus\,\mathrm d\varphi \]
का एक कदम $f =$ $2\Omega\sin\varphi$ को $2\Omega\cos\varphi\,\mathrm d\varphi$ से बदलता है, इसलिए $\beta =$ $2\Omega\cos\varphi/a_\oplus =$ $2 \times \num{7.292e-5} \times \cos\ang{30}/(\num{6.371e6}) =$
$\qty{1.98e-11}{m^{-1}.s^{-1}}$। (ख)~अब $R/\beta = \num{e-6}/(\num{1.98e-11}) =
\qty{5.04e4}{m}$, लगभग \qty{50}{km}। (ग)~अंदरूनी भाग के लिए: $\num{30e6}/(\num{5e6} \times
1000) = \qty{6.0e-3}{m/s}$, \qty{6}{mm/s}; \omterm{def:b1:ocean-circulation:wbc}{पश्चिमी सीमा धारा}: $\num{30e6}/(\num{5.04e4} \times 1000) =$ $\qty{0.60}{m/s}$, सौ गुना तेज़, चौड़ाइयों का अनुपात। खाड़ी धारा, \qty{1.8}{m/s} पर, और भी तेज़ है: वास्तविक धारा इस आकलन से संकरी है, उसका प्रवाह सतह के निकट संकेंद्रित है, और उसका पुनःपरिसंचरण उसके परिवहन में जोड़ता है।
\end{solution}

\begin{solution}{exo:b1:ocean-circulation:12}
(क)~जहाँ भी समुद्र \qty{1000}{m} से गहरा है, \qty{1000}{m} के ऊपर का स्तंभ अधिकतम \qty{1}{km} जल रखता है, और अन्यत्र कम: \qty{1000}{m} के ऊपर का आयतन अधिकतम $\num{361.84e6} \times 1 = \qty{3.618e8}{km^3}$ है, इसलिए नीचे का आयतन कम से कम $\num{1.3324e9} - \num{3.618e8} = \qty{9.71e8}{km^3}$ है। (ख)~$\num{26e6} \times \num{3.156e7} = \qty{8.21e14}{m^3}$ प्रति वर्ष, \qty{8.21e5}{km^3}; $\num{9.71e8}/(\num{8.21e5}) \approx 1200$ वर्ष, संवाहक के हज़ार वर्षों की कोटि के समान। (ग)~$(1/2)^{25/12.3} = 0.244$, लगभग एक-चौथाई; $\num{e-3} \times 25 \times \num{3.156e7} = \qty{7.9e5}{m}$, लगभग \qty{790}{km}। 25 वर्षों के बाद ट्रिटियम उतना ही फैला है जितना एक धीमा अपवाह उसे ले जाता, जबकि उसका एक-चौथाई पता लगाने योग्य बचा है।
\end{solution}

\begin{solution}{pb:b1:ocean-circulation:1}
\begin{enumerate}
  \item एक उत्तर से बहती पवन तट के साथ दक्षिण की ओर बहती है।
  \item $1000\,C_D = 0.60 + 0.071 \times 8 = 1.168$: $C_D = \num{1.168e-3}$।
  \item $T = 1.3 \times \num{1.168e-3} \times 64 = \qty{0.0972}{Pa}$।
  \item $C_D = \num{1.452e-3}$ और $T = \qty{0.272}{Pa}$: प्रतिबल $2.80$ गुना बढ़ता है,
        जो $(12/8)^2 = 2.25$ से अधिक है, क्योंकि $C_D$ स्वयं पवन के साथ बढ़ता है।
  \item $\Omega = \qty{7.292e-5}{rad.s^{-1}}$; $f = 2\Omega\sin\ang{38} =
        \qty{8.979e-5}{s^{-1}}$।
  \item $2\pi/f = \qty{7.00e4}{s}$, \qty{19.4}{h}। कोरिऑलिस बल इस समय में एक धारा को एक पूरे वृत्त में मोड़ देता है; उससे पहले कि पवन लगभग उतनी देर बही हो, जल स्थायी संतुलन में बसने के बजाय दोलन करता है।
  \item \ang{90} दक्षिण की ओर बहती पवन के दाईं ओर: पश्चिम की ओर, खुले समुद्र की ओर।
  \item $M = T/f = 0.0972/(\num{8.979e-5}) = \qty{1082}{kg.m^{-1}.s^{-1}}$।
  \item $1082/1027 = \qty{1.054}{m^2/s}$ तट के प्रति मीटर।
  \item $D_E = 7.6 \times 8/\sqrt{\sin\ang{38}} = \qty{77.5}{m}$।
  \item $1.054/77.5 = \qty{0.0136}{m/s}$, \qty{1.36}{cm/s}।
  \item $V_0 = 0.0127 \times 8/\sqrt{\sin\ang{38}} = \qty{0.129}{m/s}$, पवन के दाईं ओर \ang{45} पर: दक्षिण-पश्चिम की ओर। ड्रिफ़्टर एकमान के मान का लगभग आधा देते हैं, लगभग \qty{0.07}{m/s}।
  \item तट से परत में बहा ले जाया गया जल न तो तट के आर-पार बदला जा सकता है, न ऊपर से: वह परत के नीचे से तट के पास ऊपर उठते जल से बदला जाता है, जैसा \cref{sec:b1:ocean-circulation:upwelling} में खींचे गए परिच्छेद में है; समतापी रेखाएँ तट की ओर ऊपर मुड़ती हैं।
  \item $1.054 \times 1000 \times 86400 = \qty{9.11e7}{m^3}$ प्रति दिन प्रति किलोमीटर तट।
  \item $w = 1.054/(\num{2e4}) = \qty{5.27e-5}{m/s}$, \qty{4.55}{m} प्रति दिन।
  \item $77.5/4.55 = 17$ दिन।
  \item वायु नीचे से ठंडे जल द्वारा ठंडी होती है: सतह के निकट की परत ठंडी और स्थिर हो जाती है, नमी कोहरे और निचले स्तरी-मेघ में संघनित हो सकती है, और संवहन न होने से तट पर कम वर्षा होती है।
  \item उत्सर्पित जल पोषक तत्वों से समृद्ध है; वे फ़ाइटोप्लवकों को पोषित करते हैं, जो ज़ूप्लवकों और मछलियों को।
  \item उत्तर की ओर बहती पवन उसके दाईं ओर \ang{90} पर, पूर्व की ओर, एक परिवहन चलाती है: सतही जल तट के विरुद्ध जमा होता है और डूबता है (\omterm{def:b1:ocean-circulation:upwelling}{अवसर्पण}); \omterm{def:b1:ocean-circulation:upwelling}{उत्सर्पण} रुक जाता है और गरम सतही जल तट पर लौटता है।
  \item $f(\ang{20}) = \qty{4.988e-5}{s^{-1}}$: $M = \qty{1948}{kg.m^{-1}.s^{-1}}$,
        $\sin\ang{38}/\sin\ang{20} = 1.80$ गुना अधिक।
  \item दक्षिणी गोलार्ध में परिवहन पवन के बाईं ओर \ang{90} पर है। उसे पश्चिम की ओर, अपतटीय, धकेलने के लिए पवन को उत्तर की ओर, दक्षिण से बहना होगा: दोनों स्थितियों में पवन तट के साथ भूमध्य रेखा की ओर बहती है।
  \item $1000\,C_D = 0.29 + 0.62 + 0.308 = 1.218$; $T = \qty{0.0396}{Pa}$; $M =
        \qty{441}{kg.m^{-1}.s^{-1}}$; दोनों के लिए अनुपात $0.407$ है।
  \item तट के प्रति मीटर, $M c_p\,\Delta T = 1082 \times 3992 \times 6 =
        \qty{2.59e7}{W}$, \qty{25.9}{MW}; \qty{1000}{km} के साथ, \qty{2.59e13}{W}, \qty{25.9}{TW}।
  \item $\num{31.7e6}/1.054 = \qty{3.01e7}{m}$, लगभग \qty{30000}{km}: \omterm{def:b1:ocean-circulation:upwelling}{तटीय उत्सर्पण} \omterm{def:b1:ocean-circulation:wbc}{पश्चिमी सीमा धारा} से कहीं कम जल चलाता है।
  \item $Q = 1.054 \times \num{e6} = \qty{1.05e6}{m^3/s}$: \qty{1000}{km} तट के साथ \omterm{def:b1:ocean-circulation:upwelling}{उत्सर्पण} परिवहन \qty{1.05}{\sverdrup} है।
\end{enumerate}
\end{solution}
